\section{Multiple Qubit Gates}\label{sec:multiple-qubit-gates}

\subsection{Multiple Qubit States}\label{sec:multiple-qubit-states}

Now that we have a solid understanding of single qubits and their states, we can transition to discussing many qubits as one joint system, rather than just one vector with one qubit.

\highspace
With \textbf{one qubit}, life is simple because the state of the system can be described by a single vector in a two-dimensional complex vector space:
\begin{equation*}
    \ket{v} = a \ket{0} + b \ket{1}
\end{equation*}
With \textbf{two qubits}, something radical happens. We can no longer describe the system by looking at each qubit separately. Instead, we must describe \textbf{one global state} that captures the behavior of both qubits together. For example:
\begin{equation*}
    \ket{v} = c_{0} \ket{00} + c_{1} \ket{01} + c_{2} \ket{10} + c_{3} \ket{11}
\end{equation*}
This is not just a simple combination of two single-qubit states.

\highspace
\textcolor{Red2}{\faIcon{exclamation-triangle} \textbf{Exponential Explosion.}} The first thing to notice is that the number of amplitudes needed to describe the state of a system grows exponentially with the number of qubits.
\begin{itemize}
    \item For 1 qubit, we need 2 amplitudes: $a$ and $b$.
    \item For 2 qubits, we need 4 amplitudes: $c_{0}$, $c_{1}$, $c_{2}$, and $c_{3}$.
    \item For 3 qubits, we need 8 amplitudes: $d_{0}$, $d_{1}$, $d_{2}$, $d_{3}$, $d_{4}$, $d_{5}$, $d_{6}$, and $d_{7}$.
    \item \hl{For $n$ qubits, we need $2^{n}$ amplitudes to fully describe the state of the system.}
\end{itemize}
This is why quantum computing is so powerful: it can \textbf{process many possibilities at the same time} and \textbf{manipulate them so that only the correct ones survive}.

\highspace
\begin{flushleft}
    \textcolor{Green3}{\faIcon{book} \textbf{The Key New Phenomenon: Entanglement}}
\end{flushleft}
\textbf{Entanglement is the key new phenomenon when dealing with multiple qubits.} It refers to quantum states that cannot be expressed as a tensor product of individual qubit states. For example, the state:
\begin{equation*}
    \dfrac{1}{\sqrt{2}} \left( \ket{00} + \ket{11} \right)
\end{equation*}
Is an entangled state. In such cases, the system must be described globally, and it is not possible to assign independent quantum states to each qubit.

\highspace
If one qubit is measured, the outcome is random, but it determines the outcome of the other qubit due to the structure of the global state. This leads to strong correlations between measurements results. However, this does not imply any form of information transfer between the qubits.

\highspace
Entanglement is a fundamental resource in quantum computing, enabling correlations that have no classical counterpart, especially when combined with interference effects in quantum algorithms.
\begin{itemize}
    \item Entanglement does not mean that one qubit sends information to the other. The observed correlations arise from the collapse of a \textbf{global quantum state}, not from communication between subsystems.
    
    \item Before measurement, individual qubits in an entangled state do \textbf{not have a well-defined state}. Only the joint system is well-defined.
    
    \item Each qubit individually appears \textbf{random}, but measurement outcomes are \textbf{correlated}.
    
    \item Entanglement alone is not the source of quantum computational power; it becomes powerful when combined with \textbf{interference} and quantum operations.
\end{itemize}

\highspace
\begin{flushleft}
    \textcolor{Green3}{\faIcon{atom} \textbf{New tools we need}}
\end{flushleft}
To handle this new world of multiple qubits, we need new mathematical tools. We will introduce the \textbf{tensor product}, which allows us to construct the joint state of multiple qubits; we will also introduce the \textbf{multi-qubit gates}, which are operations that can manipulate the joint state of multiple qubits simultaneously; and we will discuss how to \textbf{measure} multiple qubits and understand the resulting probabilities.